
\begingroup
\setcounter{proposition}{11}% Conservar el número visible XIII.2.12 de la fuente.
\begin{corollaire}
\label{XIII.2.12}
Sea $X$ una curva propia, lisa y conexa de género $g$ sobre un cuerpo
separablemente cerrado $k$ de característica $p\geq 0$. Sea $U$ el abierto
obtenido al quitar de~$X$ $n$ puntos cerrados distintos
$a_1,\dots,a_n$. Entonces el grupo fundamental moderadamente ramificado
$\pi_1^{\tame}(U)$ \eqref{XIII.2.1.3} puede ser engendrado por $2g+n$
elementos $x_i,y_i,\sigma_j$, con $1\leq i\leq g$ y
$1\leq j\leq n$, de modo que $\sigma_j$ sea un generador de un grupo
de inercia correspondiente a~$a_j$, y se tenga la relación
\begin{equation*}
\label{eq:XIII.2.12.*}
\tag{$*$} \prod_{1\leq i\leq
g}(x_iy_ix^{-1}_iy^{-1}_i)\cdot\prod_{1\leq j\leq n}\sigma_j=1.
\end{equation*}
Para todo grupo finito $G$ \emph{de orden primo con} $p$,
engendrado por elementos $\overline{x}_i,\overline{y}_i,
\overline{\sigma}_j$ que satisfacen la relación
\eqref{eq:XIII.2.12.*}, existe un recubrimiento étale de~$U$, de grupo
$G$, correspondiente a un homomorfismo $\pi_1^{\tame}(U)\to G$ que
\marginpar{393}%
envía $x_i,y_i,\sigma_j$ a~$\overline{x}_i,\overline{y}_i,
\overline{\sigma}_j$, respectivamente. En otros términos, si $p'$ denota
el conjunto de los números primos distintos de~$p$, $\pi_1^{p'}(U)$ es el
grupo pro-$p'$ engendrado por los generadores $x_i,y_i,\sigma_j$, sujetos
a la única relación~\eqref{eq:XIII.2.12.*}.
\end{corollaire}

\subsubsection*{Demostración}
Podemos suponer que $k$ es algebraicamente cerrado. Supongamos primero que
$k$ tiene característica cero. Existe entonces una subextensión
algebraicamente cerrada $k'$ de~$k$, de grado de trascendencia finito
sobre~$\QQ$, tal que $X$ provenga por extensión de escalares de una curva
propia y lisa $X'$ definida sobre~$k'$, y podemos suponer que los puntos
$a_1,\dots,a_n$ provienen de puntos racionales $a'_1,\dots,a'_n$ de~$X'$.
Como $k'$ tiene grado de trascendencia finito sobre~$\QQ$, se puede hallar
una inmersión de~$k'$ en el cuerpo de los números complejos $\CC$; sea
$\widetilde{U}=U'\times_{k'}\CC$. Sea $k''$ una extensión algebraicamente
cerrada de~$k'$ tal que existan morfismos sobre~$k'$ de~$k$ y de~$\CC$
en~$k''$. Si $g'\colon U'\to k'$ es el morfismo estructural y
\ifthenelse{\boolean{orig}}
{si $\cal F$ es un haz de grupos constante finito sobre~$U''$}
{si $\cal F'$ es un haz de grupos constante finito sobre~$U'$},
se deduce de~\Ref{XIII.2.9} que los morfismos de especialización
$$
(\R^1g'_*\cal{F}')_{\CC}\to(\R^1g'_*\cal{F}')_{k''}
\from(\R^1g'_*\cal{F}')_k
$$
son isomorfismos. En términos de grupos fundamentales, esto demuestra que
se tiene un isomorfismo, definido salvo automorfismo interior,
$$
\pi_1(U)\to\pi_1(\widetilde{U}),
$$
y está claro que este isomorfismo transforma un grupo de inercia relativo
a un punto de~$X'-U'$ en un grupo de inercia relativo al mismo punto.
Por tanto, podemos suponer que $k=\CC$. En este último caso, del teorema
de existencia de Riemann (\Ref{XII.5.2}) resulta que el grupo fundamental
$\pi_1(U)$ no es sino la completación, para la topología de los subgrupos
de índice finito, del grupo fundamental del espacio analítico asociado
a~$U$. Ahora bien, este último puede calcularse por métodos trascendentes
\cite[ch\ptbl 7
\S 47]{XIII.3};
\marginpar{394}%
puede ser engendrado por $2g+n$ elementos $x_i,y_i,\sigma_j$ tales que
$\sigma_j$ sea la imagen de un generador
\ifthenelse{\boolean{orig}}
{del grupo fundamental local $\pi_1(D_j)$ de un pequeño disco centrado
en~$a_j$}
{del grupo fundamental local $\pi_1(D_j-\{a_j\})$, donde $D_j$ es un
pequeño disco centrado en~$a_j$},
\ie un generador de un grupo de inercia correspondiente al punto $a_j$,
y estos elementos satisfacen la única relación \eqref{eq:XIII.2.12.*}.

Si ahora $k$ tiene característica $p>0$, se puede hallar un anillo de
valoración discreta completo $A$, con cuerpo residual $k$ y cuerpo de
fracciones $K$ de característica cero, y un esquema conexo $X_1$, propio
y liso sobre~$S=\Spec A$, tal que
$X_1\times_S\Spec k\simeq X$ (\Ref{III.7.4}). Los puntos $a_j$ se
levantan entonces a secciones $s_j$ de~$X_1$ sobre~$S$; sea $Y_{1j}$ el
subesquema cerrado reducido cuyo espacio subyacente es $s_j(S)$, sea $Y_1$
la unión de los $Y_{1j}$, sea $U_1=X_1-Y_1$ y sea
$g_1\colon U_1\to S$ el morfismo estructural. Sean $\overline K$ una
extensión algebraicamente cerrada de~$K$,
\ifthenelse{\boolean{orig}}
{}
{$\overline X=X_1\times_S\overline K$ y}
$\overline U=U_1\times_S\overline K$. Si $C$ es un grupo constante finito,
se deduce de~\Ref{XIII.2.8} que el morfismo de especialización
$$
(\R^1g_{1*}C_{U_1})_k\to(\R^1g_{1*}C_{U_1})_{\overline K}
$$
es inyectivo, e incluso biyectivo si $C$ es de orden primo con~$p$.
Ahora bien, esto significa, en términos de grupos fundamentales, que el
morfismo de especialización
\ifthenelse{\boolean{orig}}
{\eqref{XIII.1.10}}
{\eqref{XIII.2.10}}
$$
\pi\colon \pi_1(\overline U)\to\pi^{\tame}_1(U)
$$
es sobreyectivo, y que el morfismo de especialización
$$
\pi^{p'}_1(\overline U)\to\pi^{p'}_1(U)
$$
es biyectivo. Por último, si $x_i,y_i,\sigma_j$ son generadores de
$\pi_1(\overline U)$ tales que $\sigma_j$ sea un generador de un grupo de
inercia correspondiente al punto $b_j=Y_{1j}(\overline K)$ de~$\overline X$,
entonces, según
\ifthenelse{\boolean{orig}}
{\Ref{XIII.1.11}}
{\Ref{XIII.2.11}},
$\pi(\sigma_j)$ es un generador de un grupo de inercia correspondiente
a~$a_j$, lo cual concluye la demostración.

\endgroup
